%-----------------------------------------------------------------------
\section{Limitations and Future Work}
%-----------------------------------------------------------------------

\paragraph{Benchmark scope.}
Our evaluation processes raw frames through shuttle tracking, pose extraction,
feature construction, and stroke classification. However, stroke windows still
use ground-truth hit-frame indices and court normalization relies on provided
homographies. The system should therefore be viewed as a raw-feature and
stroke-classification pipeline under supplied temporal and spatial metadata, not
a fully automatic system for arbitrary unlabeled broadcasts.

\paragraph{Perception quality.}
The main error source is far-side pose quality: its validity is just 54.29\%.
Rally filtering occasionally retains post-rally broadcast content, adding noise
near rally endpoints. Future work should focus on improving far-side detection,
tightening rally boundaries, and adopting trajectory-based hit detection to
remove reliance on annotated hit frames.

\paragraph{Output scope.}
We do not yet predict point winners, shuttle landing positions, persistent
player identities, or continuous footwork trajectories. Side-specific
distributions and clusters are therefore descriptive rather than causal: Top
and Bottom are camera-relative positions, not persistent athletes, so any side
differences reflect court geometry rather than player identity. Extending the
system to predict outcomes and integrate point-level context is a promising
direction.

\paragraph{Future work.}
Future extensions should improve far-side pose detection, replace ground-truth
hit frames with trajectory-based hit detection, automate court calibration, and
refine fine-grained stroke classification. Integrating the system with VR-based
visual analytics such as VIRD \cite{vird2023} could also provide practitioners
with interactive exploration of automatically derived event streams.